\subsection{Several crossing base divisors}
\label{sec:crossing-fiber-integrability}

In \cref{thm:fiber-integrability-band} the upper bounds are required only on one polydisk over each divisor, which can be chosen away from the crossings; integrability at the crossings is controlled by the global lower bound. Exact models can nevertheless have fiber poles over the crossings, as in the following example, where different divisors use separate fiber coordinates.

\Needspace{32\baselineskip}
\begin{example}\label{ex:crossing-fiber-sharp-error}\index{Holomorphic section!crossing sharp error}
Let $B=(\mathbb P^1)^q$, $X=B\times(\mathbb P^1)^q$ and $A=\mathcal O_B(1,\ldots,1)$. For each $i$, let $D_i\subset B$ be the inverse image of a marked point in the $i$th factor of $B$, and let $E_i\subset X$ be the inverse image of a marked point in the $i$th factor of the fiber. Use normalized Fubini--Study metrics on the corresponding line bundles, and write $\omega_{B,i},\omega_{F,i}$ for the pullbacks to $X$ of the normalized Fubini--Study forms on the $i$th factors of $B$ and of the fiber. Put
\[
 \begin{gathered}
 R_i=\|\pi^*s_i\|^2,\quad T_i=\|s_{E_i}\|^2,\quad
 \varphi=\frac14\sum_{i=1}^q\log(R_i^2+T_i^3),\\
 \omega=\sum_i(\omega_{B,i}+\omega_{F,i}),\qquad
 \psi=\frac12\sum_i\log\|s_i\|^2.
 \end{gathered}
\]
Then $\varphi\in\PSH(X,\omega)$, and \cref{thm:fiber-integrability-band} applies with $N=q$, $u=0$, $c_i=1/2$, $p_i=2$, $b_{i1}=3$ and $a_i=2/3$. The reference potential $\psi$ has mass $q!/2^q$. With $L_m=\lfloor m/2\rfloor$ and $K_m=\max\{\lfloor m/2-2/3\rfloor,0\}$, the selected and reference spaces are exactly
\[
 \begin{split}
 V_m&=H^0\bigl(B,\mathcal O_B(m-K_m,\ldots,m-K_m)\bigr),\\
 W_m&=H^0\bigl(B,\mathcal O_B(m-L_m,\ldots,m-L_m)\bigr).
 \end{split}
\]
They agree for $m=1$, and for $m\geq2$ the defect is
\[
 \dim V_m-\dim W_m
       =(\lceil m/2\rceil+2)^q-(\lceil m/2\rceil+1)^q.
\]
For $q\geq2$ this is $q2^{1-q}m^{q-1}+O(m^{q-2})$, and for $q=1$ it is exactly one.
\end{example}
\begin{proof}
By the calculation in \cref{ex:anisotropic-fiber-sharp-error}, the complex Hessian of each summand of $\varphi$ is bounded below by $-\tfrac12\omega_{B,i}-\tfrac34\omega_{F,i}$, also across its pole set by the extension theorem after a quadratic correction. Summing proves $\varphi\in\PSH(X,\omega)$. The global lower bound follows from $R_i^2+T_i^3\geq R_i^2$. For the upper bound over $D_i$, away from the other divisors, place the $i$th fiber coordinate at its marked point; all the other summands are smooth there. The residual is $\beta=\tfrac12\sum_i\omega_{B,i}>0$, whose top integral is $q!/2^q$.

The theorem gives the required generic orders along every $D_i$. Conversely, near a point where several of the pairs $(R_i,T_i)$ vanish, the density $\mathrm{e}^{-m\varphi}$ is comparable to the product of the factors $(|z_i|^4+|w_i|^6)^{-m/4}$ in disjoint pairs of variables. A descended coefficient divisible by every $z_i^{K_m}$ has a bounded holomorphic residual. Each factor is integrable because $K_m+1>m/2-2/3$, and Fubini's theorem gives integrability at the intersections; elsewhere the same argument uses fewer factors. This proves the identity for $V_m$, and the theorem gives $W_m$. A line bundle $\mathcal O_B(d,\ldots,d)$ with $d\geq0$ has $(d+1)^q$ independent sections, the products of monomials in the separate factors. Since $K_m=L_m-1$ for $m\geq2$ and $m-L_m=\lceil m/2\rceil$, subtraction gives the displayed defect and its leading term.
\end{proof}

As before, the selecting potential takes finite values and the value $-\infty$ in the fibers over the divisors, so it does not descend.

\subsection{Fiber integrability on a modified base}
\label{sec:modified-base-fiber-integrability}

Retain $\pi:X\to B$, $(A,h)$, $\theta$, $\varphi$ and $V_m$ from \cref{sec:fiber-integrability-band}. A modification of the base can turn a higher-codimension singularity into divisorial data. The integrability test then acts on the pulled-back section multiplied by the Jacobian. This fixed factor changes the exact orders, while the restriction estimate still bounds the section loss.

Let $\rho:Y\to B$ be a modification of connected compact complex manifolds. Form the fiber product and its projections
\[
 X'=X\times_B Y,\qquad \pi':X'\to Y,\qquad \mu:X'\to X.
\]
Local submersion coordinates show that $X'$ is smooth, $\pi'$ is a submersion with the same compact connected fibers, and $\mu$ is a modification. The relative canonical bundle $R_\rho$ and Jacobian section $J_\rho$ are those of \cref{sec:psh-birational-ideals}. The coordinate Jacobian of $\mu$ is the base Jacobian times the identity in the fiber variables, giving
\[
 R_\mu=\pi'^*R_\rho,\qquad J_\mu=\pi'^*J_\rho.
\]
Put $\varphi'=\varphi\circ\mu$ and $\theta'=\rho^*\theta$. Choose distinct nonempty smooth connected reduced divisors $D_1,\ldots,D_N$ on $Y$ with simple-normal-crossing union, smooth metrics $H_i$, defining sections $s_i$ and curvatures $\eta_i$.

\Needspace{14\baselineskip}
\begin{theorem}\label{thm:modified-base-fiber-band}\index{Holomorphic section!modified base band}
Suppose $c_i>0$ and $u\in C^\infty(Y,\mathbb R)$ satisfy
\[
 \beta=\theta'+\ddc u-\sum_i c_i\eta_i>0,
 \qquad \varphi'\geq\sum_i c_i\log\|\pi'^*s_i\|_{H_i}^2-C.
\]
For each $D_i$, assume one anisotropic upper witness for $\varphi'$ and $\pi'$ away from the other divisors, as in \cref{thm:crossing-fiber-integrability-band}, with fixed $p_i,b_{ij}$. Define $a_i,L_{m,i},K_{m,i}$ by the formulas in that theorem, and put
\[
 \psi'=u+\sum_i c_i\log\|s_i\|_{H_i}^2,
 \qquad M_m=\rho^*A^m\otimes R_\rho.
\]
There is a unique $\psi\in\PSH(B,\theta)$ with $\psi\circ\rho=\psi'$, and
\[
 \mathrm{MA}_\theta(\psi)=\rho_*(\beta^q),\qquad
 \int_B\mathrm{MA}_\theta(\psi)=\int_Y\beta^q\in(0,\infty).
\]
The map $T_m(s)=J_\rho\rho^*s$ identifies $H^0(B,A^m)$ with $H^0(Y,M_m)$. For $W_m=H^0(B,A^m\otimes\mathcal I(m\psi))$,
\[
 \begin{split}
 T_m(W_m)&=H^0\left(Y,M_m\left(-\sum_iL_{m,i}D_i\right)\right),\\
 T_m(W_m)&\subseteq T_m(V_m)
       \subseteq H^0\left(Y,M_m\left(-\sum_iK_{m,i}D_i\right)\right).
 \end{split}
\]
Consequently $W_m\subseteq V_m$ and $\dim V_m-\dim W_m=O(m^{q-1})$. On a curve the defect is at most $\sum_i\lceil a_i\rceil$ for every $m$.
\end{theorem}
\begin{proof}
Poincar\'e--Lelong gives $\theta'+\ddc\psi'=\beta+\sum_i c_i[D_i]$. Thus $\psi'\in\PSH(Y,\rho^*\theta)$, and \cref{prop:psh-descent-modification} gives the unique descended potential. The non-pluripolar measure upstairs is $\beta^q$ by locality off the divisor union and exclusion of that union. Apply the measure identity in \cref{prop:birational-partial-envelopes} to obtain the asserted pushforward and mass.

\Cref{lem:adjoint-ideal-transform}, first with weight zero and then with $m\psi$, gives the whole-space identification and the identity for $T_m(W_m)$. The simple-normal-crossing integrability calculation in \cref{thm:crossing-fiber-integrability-band} supplies $\mathcal I(m\psi')=\mathcal O_Y(-\sum_iL_{m,i}D_i)$. Boundedness of $u$ and the global lower bound give $\varphi'\geq\psi'\circ\pi'-C'$. Since $\rho\circ\pi'=\pi\circ\mu$ and $\mu$ is surjective, this descends to $\varphi\geq\psi\circ\pi-C'$. Integrability in bounded submersion charts proves $W_m\subseteq V_m$.

For $s\in V_m$, the adjoint transformation for $\mu$ gives
\[
 J_\mu\mu^*s=\pi'^*T_m(s)
 \in H^0\bigl(X',\pi'^*M_m\otimes\mathcal I(m\varphi')\bigr).
\]
Its descended coefficient in each upper witness therefore has order at least $K_{m,i}$, by \cref{lem:anisotropic-fiber-integrability} and the generic unit-coefficient argument of \cref{thm:fiber-integrability-band}. This order is the order of $T_m(s)$, including its Jacobian factor. It proves the upper inclusion.

The divisor restriction chain on $Y$ has at most $\sum_i\lceil a_i\rceil$ steps. On $D_i$, each target twists $(\rho^*A|_{D_i})^m$ by the powers $\mathcal O_Y(-D_j)|_{D_i}$ in a linear range and by the fixed bundle $R_\rho|_{D_i}$. Include this last bundle with exponent one in \cref{lem:finite-family-section-bound}; every target has dimension $O(m^{q-1})$ uniformly. Summing gives the estimate. For $q=1$, each divisor is a point and every target has dimension one.
\end{proof}

The theorem uses a specified smooth modification and divisorial data with a positive residual. It does not assume that arbitrary selecting potentials furnish these data. The following point-pole model makes both the residual and the Jacobian correction explicit.

\Needspace{14\baselineskip}
\begin{example}\label{ex:modified-point-fiber-series}\index{Holomorphic section!point pole Jacobian}
Let $q\geq2$, $B=\mathbb P^q$, $A=\mathcal O_B(1)$, $p=[1:0:\cdots:0]$ and $X=B\times\mathbb P^1$. In normalized Fubini--Study metrics, for a marked fiber point $o$ put
\[
 R=\frac{\sum_{j=1}^q|Z_j|^2}{\sum_{j=0}^q|Z_j|^2},\quad
 T=\|s_o\|^2,\quad \varphi=\tfrac14\log(R^2+T^3),\quad
 \psi=\tfrac12\log R.
\]
Then $\varphi\in\PSH(X,\omega_B+\omega_F)$ and $\psi\in\PSH(B,\omega_B)$ has mass $1-2^{-q}$. With the ideal sheaf $\mathfrak m_p$ of $p$, the exact spaces are
\[
 \begin{gathered}
 V_m=H^0\bigl(B,A^m\otimes\mathfrak m_p^{k_m}\bigr),\qquad
 W_m=H^0\bigl(B,A^m\otimes\mathfrak m_p^{\ell_m}\bigr),\\
 k_m=\max\{\lfloor m/2-2/3\rfloor-q+1,0\},\qquad
 \ell_m=\max\{\lfloor m/2\rfloor-q+1,0\}.
 \end{gathered}
\]
Their defect vanishes for $m<2q$. For $m\geq2q$ it is
\[
 \dim V_m-\dim W_m=\binom{\lfloor m/2\rfloor-1}{q-1}
       \sim\frac{2^{1-q}}{(q-1)!}m^{q-1}.
\]
\end{example}
\begin{proof}
A common denominator writes $R^2+T^3$ as a sum of squared holomorphic section coefficients in $\mathcal O_X(2,3)$ divided by its smooth metric denominator. Hence $\ddc\varphi\geq-\tfrac12\omega_B-\tfrac34\omega_F$, including across the common zero set by local psh extension.

Let $\rho:Y\to B$ be the blowup of $p$, with exceptional divisor $E$, and let $\tau:Y\to\mathbb P^{q-1}$ be the resolved projection $[Z_0:\cdots:Z_q]\mapsto[Z_1:\cdots:Z_q]$. The graph embedding $(\rho,\tau)$ shows that
\[
 \alpha=\rho^*\omega_B,\quad\gamma=\tau^*\omega_{\mathbb P^{q-1}},\qquad
 \beta=\tfrac12(\alpha+\gamma)>0.
\]
The divisor bundle is $\mathcal O_Y(E)=\rho^*\mathcal O_B(1)\otimes\tau^*\mathcal O_{\mathbb P^{q-1}}(-1)$. Its quotient metric has curvature $\alpha-\gamma$ and satisfies $R\circ\rho=\|s_E\|^2$. Thus the theorem applies with $u=0$, $c=1/2$, $a=2/3$, and residual $\beta$. The pulled-back selector has the exact model $\tfrac14\log(|z|^4+|w|^6)+O(1)$ along $E\times\{o\}$ and is locally bounded elsewhere. Divisibility by $z^{K_m}$, where $K_m=\max\{\lfloor m/2-2/3\rfloor,0\}$, gives the reverse upper inclusion by \cref{lem:anisotropic-fiber-integrability} in every pole chart.

In blowup coordinates $Z_1/Z_0=z$, $Z_j/Z_0=zt_j$ for $j\geq2$, the Jacobian is $z^{q-1}$. A section of base order $n$ therefore gives $T_m(s)$ of generic order $n+q-1$ along $E$. The two divisorial thresholds yield exactly $k_m,\ell_m$ above. Homogeneous monomials realize every jet of total degree at most $m$ at $p$. For $m\geq2q$, $k_m=\ell_m-1$ and the newly allowed homogeneous degree $\ell_m-1$ has $\binom{\ell_m+q-2}{q-1}$ monomials, giving the defect. For $m<2q$ both orders are zero.

Finally, in cohomology $[\beta]=\rho^*c_1(A)-\tfrac12[E]$. The mixed intersection numbers containing both $\rho^*c_1(A)$ and $[E]$ vanish because the pullback restricts to zero on $E$. The normal bundle of $E\simeq\mathbb P^{q-1}$ is $\mathcal O_E(-1)$, so $\int_Y[E]^q=(-1)^{q-1}$, while $\int_Y(\rho^*c_1(A))^q=1$. Expanding gives $\int_Y\beta^q=1-2^{-q}$, which is the mass of $\psi$ by the theorem.
\end{proof}

The selector is finite at most points of the fiber over $p$ and equals $-\infty$ at $o$, so it does not descend up to bounded difference. Its base reference has a higher-codimension pole; the Jacobian is what converts the resolved divisor orders into the base point-ideal orders.

